\section{Limitações e Ameaças à Validade}

O trabalho apresenta limitações em quatro frentes. Primeiro, a cobertura da coleta depende da infraestrutura digital das câmaras municipais. Municípios com melhor padronização tecnológica tendem a ficar sobre-representados, introduzindo viés regional e institucional. A análise exploratória mostrou, por exemplo, que Sul e Sudeste concentram 68,55\% dos registros e 64,29\% dos municípios do dataset final. Segundo, a avaliação do tradutor ainda carece de \textit{baselines} mais fortes e de um conjunto de teste independente, já que o BERTScore reportado foi medido no conjunto de validação. Terceiro, a busca semântica ainda não dispõe de benchmark anotado nem métricas formais de recuperação, o que limita a força das conclusões sobre C3. Quarto, a análise de indicadores é observacional e não causal: outros fatores públicos, administrativos e socioeconômicos podem explicar parte das variações observadas.

Também há limitações específicas da proxy educacional baseada em matrículas e do agrupamento textual por Jaccard. Embora úteis operacionalmente, essas escolhas simplificam fenômenos complexos e devem ser interpretadas como parte de uma avaliação inicial de sistema, e não como evidência causal definitiva sobre efetividade de políticas públicas. Por fim, a própria base final demanda saneamento adicional: foram encontrados poucos outliers temporais importantes e cerca de 3,54\% de registros que não correspondem a PLs estritos, mas sim a anteprojetos, emendas, substitutivos ou vetos recuperados pelo filtro textual de \textit{tipo\_label}.
